\section{Conclusion}
\label{sec:conclusion}

Modern intelligent systems increasingly integrate retrieval systems, workflow engines, persistent memory, reasoning models, and human feedback. Although these components have individually advanced, they frequently maintain separate representations of system state, limiting their ability to continually adapt through accumulated operational experience.

This paper introduced the Semut Adaptive Architecture, a unified architectural framework organized around three complementary memory structures: Artifact Memory, Action Graph, and Feedback Memory. Rather than treating these components as isolated subsystems, the proposed architecture integrates them through a continuous adaptive feedback loop in which every execution contributes additional contextual information for future reasoning.

A distinguishing characteristic of the proposed architecture is its independence from any particular reasoning engine. Rule-based systems, Large Language Models, future learning algorithms, and human operators may all utilize the same adaptive memory while remaining free to evolve independently. By separating adaptive memory from reasoning, the architecture seeks to provide a stable foundation capable of supporting heterogeneous intelligent systems over time.

The primary contribution of this work is therefore architectural rather than algorithmic. Instead of introducing a new reasoning model, the proposed framework provides a common abstraction through which artifacts, executable actions, and accumulated feedback may be represented as complementary forms of adaptive memory. This unified representation offers a conceptual foundation for intelligent assistants, enterprise automation, software engineering agents, robotics, and future multi-agent systems.

As intelligent systems continue to expand across increasingly diverse software environments, architectures capable of preserving long-term experience while remaining independent of individual reasoning technologies may become increasingly valuable. The Semut Adaptive Architecture represents one possible step toward such adaptive systems by organizing memory, execution, and feedback within a coherent and extensible architectural framework.
